

\Needspace{95mm}
\TypeQuestion{}{angle_bisector_theorem_derivation}
\FigProblemBlock{47mm}{fig/lesson_06_bisector_derivation}{
\Text{$\Tri{ABC}$ の $\Ang{BAC}$ の二等分線と辺 $\Lseg{BC}$ の交点を $\mathrm{D}$ とする。点 $\mathrm{C}$ を通り $\Lseg{AD}$ に平行な直線と、辺 $\Lseg{BA}$ を点 $\mathrm{A}$ の側に延長した直線との交点を $\mathrm{E}$ とする。次の問いに答えなさい。}
\QQ{$\Ang{ACE}=\Ang{CEA}$ を示し、$\Lseg{AE}=\Lseg{AC}$ となることを説明しなさい。}{}
\QQ{$\Tri{BDA}\Similar\Tri{BCE}$ となることを証明しなさい。}{}
\QQ{前の2つの小問の結果を用いて、$\Lseg{BD}:\Lseg{BC}=\Lseg{AB}:(\Lseg{AB}+\Lseg{AC})$ となることを示しなさい。}{}
\QQ{前の小問の比例式から、$\Lseg{BD}:\Lseg{DC}=\Lseg{AB}:\Lseg{AC}$ となることを示しなさい。}{}
}{
}

\AnswerColumnBreak
\TypeQuestion{}{angle_bisector_theorem_parallel_application}
\FigProblemBlock{49mm}{fig/lesson_06_bisector_parallel}{
\Text{$\Tri{ABC}$ で、$\Lseg{AB}=6\mathrm{cm}$、$\Lseg{AC}=9\mathrm{cm}$、$\Lseg{BC}=10\mathrm{cm}$ とする。$\Ang{BAC}$ の二等分線と辺 $\Lseg{BC}$ の交点を $\mathrm{D}$ とする。点 $\mathrm{D}$ を通り $\Lseg{AC}$ に平行な直線と辺 $\Lseg{AB}$ の交点を $\mathrm{E}$ とする。次の問いに答えなさい。}
\QQ{$\Lseg{BD},\Lseg{DC}$ の長さをそれぞれ求めなさい。}{}
\QQ{$\Tri{BDE}\Similar\Tri{BCA}$ となることを証明しなさい。}{}
\QQ{$\Lseg{DE}$ の長さを求めなさい。}{}
}{
}
